\documentclass[11pt,a4paper]{article}
\usepackage[utf8]{inputenc}
\usepackage{authblk}
\usepackage{hyperref}
\usepackage{graphicx}
\usepackage{geometry}
\usepackage{booktabs}
\usepackage{float}
\usepackage{enumitem}
\usepackage{amsmath}

\geometry{margin=1in}

\title{Outlier - Measuring AI Use: A Governance Framework for Carbon, Authorship Erosion, and AI Adoption}

\author[1]{Roshan Abraham}
\affil[1]{Independent Researcher and UX Design Engineer, Ho Chi Minh City, Vietnam}

\date{22 June 2026}

\begin{document}

\maketitle

\begin{abstract}
We introduce a local-telemetry methodology to quantify the invisible externalities of AI coding assistants at the developer terminal. Although the existing literature measures the systemic risks of AI at the datacenter or organizational level, adoption of AI occurs at the individual terminal. At this intersection, two hidden costs emerge: the carbon emission of inference and the gradual dis-empowerment of the human author. We built a simple monitoring approach that tracks both. It reads the logs already sitting on the developer's machine. No new infrastructure. No data leaves the country. Why measure at the developer's computer instead of the datacenter? Because local measurement reveals a hidden geographic tax. Developers in regions with coal-heavy grids, like Vietnam, generate up to 31 times more carbon per prompt than their European counterparts. Because local measurement reveals a hidden geographic tax. Developers in regions with coal-heavy grids, such as Vietnam, generate up to 31 times more carbon per prompt than their European counterparts. Inadvertently, they pay a massive, disproportionate environmental cost just to access the same global AI tools. To prove this works, we built and deployed a live, open-source CLI instrument. We demonstrate its viability on a live repository, corroborating the high rates of AI authorship observed in public enterprise frameworks. This instrument makes both the carbon and skill costs visible at the exact moment of delegation. Why should governments care? Because AI adoption and usage monitoring gives regulators a simple, auditable signal of foreign AI dependency at three levels they already regulate: the worker (premature de-professionalization), the firm (human oversight of high-risk AI), and the nation (trade-off between foreign cloud spending and local carbon cost), all without exporting citizen data to foreign servers.
\end{abstract}

\section{Introduction}
The usual AI safety picture is sudden: a system defects and damage follows. A quieter failure is already happening. Incremental AI improvements erode human influence over the economy, culture, and the state until that loss becomes hard to reverse---not because any system defected, but because each small act of delegation is locally rational \cite{kulveit2025gradual}.

This problem is particularly acute in software engineering. A developer who asks AI to write code consumes electricity and simultaneously reduces their own engagement with the work. Existing governance frameworks treat these as separate problems. However, at the developer's computer, they are a single event. Our contributions are:
\begin{itemize}
    \item A methodology to measure AI's impact at the individual developer level using existing local logs.
    \item Two quantifiable indicators: the Authorship Ratio and the Session Carbon Footprint.
    \item An empirical case study utilizing an open-source CLI instrument demonstrating the viability of this measurement and highlighting severe geographic carbon disparities.
\end{itemize}

\section{Related Work}
\textbf{Disempowerment:} Kulveit et al. (2025) argue incremental AI substitution quietly erodes human influence \cite{kulveit2025gradual}. Recent studies show that higher Generative AI confidence predicts less critical thinking, and experienced developers become slower while believing they are faster \cite{metr2025}.

\textbf{Carbon:} The literature meters AI's footprint at model and datacenter scale: training energy and inference energy converging on $\sim$0.3 Wh per chat query \cite{oviedo2025energy, you2025energy}. We meter by developer, at the moment of delegation, weighted by the local power grid.

\textbf{Authorship:} Ecosystem trackers note fully-AI pull requests rising rapidly \cite{greptile2026}. We localize this to the individual terminal, providing an auditable proxy for human oversight.

\section{Methodology}
\subsection{Execution Share (Authorship via Line Survival)}
Past ecosystem trackers have relied on parsing \texttt{Co-Authored-By} Git trailers to measure AI involvement. However, commit-level tracking is fundamentally flawed: squashed pull requests destroy the metadata, and a commit containing one line of AI code is weighted identically to one containing ten thousand lines.

We introduce a \textbf{blame-based, content-matched line survival} metric. The metric is defined as:
\begin{equation}
Execution\_Share = \frac{Lines_{AI\_Authored\_Surviving}}{Lines_{Total\_in\_HEAD}}
\end{equation}
The local CLI instrument operates by parsing the developer's raw AI session transcripts (e.g., Claude Code logs) for \texttt{Edit} and \texttt{Write} events. It extracts the exact lines generated by the agent, hashes them, and cross-references them against \texttt{git blame} and \texttt{git ls-files} in the current \texttt{HEAD}. Crucially, this only counts lines that \textit{survive} into the living artifact. Trivial lines (blank lines, bare brackets, $<$8 characters) and generated files (e.g., lockfiles) are strictly excluded from both the numerator and denominator to prevent ratio inflation.

For agents that do not leave parseable local logs (e.g., Copilot, Aider, Cursor), we introduce an effect-observer primitive (\texttt{outlier watch}). This brackets a coding session, captures all file changes occurring within that temporal window, and attributes the surviving lines to the active tool, rendering the execution signal entirely tool-agnostic. We retain the commit-trailer proxy solely as a fallback floor when transcript data is purged or unavailable.

\subsection{Session Carbon Footprint}
For the compute side, we estimate the carbon produced by one work session using the formula:
\begin{equation}
Session\_Carbon = \sum (Tokens_{out} \times Energy_{token} \times Grid_{local\_intensity})
\end{equation}
By fixing the $Energy_{token}$ estimate ($\sim$0.3 Wh) and utilizing regional grid intensity data \cite{hust2024, rte2024}, we can isolate the geographic variable.

\section{Empirical Study \& Results}
Using our local CLI instrument, we analyzed an active 46-day solo-founder repository (160 commits) alongside public framework data.

\subsection{Authorship Findings}
Using the fallback commit-proxy on a 160-commit history, we established an initial baseline of 75.6\% AI authorship. However, when applying the rigorous blame-based execution tracking over the repository's 8.2K live lines of code, the instrument revealed a true Execution Share of 62\% (5.1K lines directly attributable to AI generation that survived into production). We verified the commit-level proxy against public repositories (e.g., Microsoft Semantic Kernel, Vercel AI SDK), which showed AI authorship shares ranging from 21\% to 79\%, confirming the tool's ability to measure both deep line-survival and broad commit proxies.

\subsection{The Geographic Carbon Tax}
Aggregating the developer's Claude Code session logs yielded 2.26 billion tokens processed, with 94.9\% representing redundant cache reads. Based on output token generation ($\sim$10 kWh inference energy), we applied grid intensity factors:
\begin{itemize}
    \item Vietnam (681 $gCO_2/kWh$): 6.81 $kgCO_2$
    \item USA (384 $gCO_2/kWh$): 3.84 $kgCO_2$
    \item France (21.7 $gCO_2/kWh$): 0.22 $kgCO_2$
\end{itemize}
This reveals an undeniable infrastructural fact: the same AI workload produces approximately 31 times more carbon on Vietnam's grid than on France's grid. 

\section{Socio-Economic \& Governance Implications}
AI Authorship Ratio and Session Carbon provide the state with a leading indicator of AI dependency. For IT-export economies in the Global South, a high AI authorship share threatens premature deprofessionalization. Simultaneously, the massive carbon disparity highlights an ``imported emissions'' problem, where local infrastructure is forced to absorb the environmental cost of accessing foreign AI tools.

This local-first telemetry allows governments to enforce mandates like Vietnam's Decree 142/2026 (requiring human oversight for high-risk AI) without exporting citizen data to foreign servers.

\section{Conclusion}
Gradual disempowerment and AI's carbon cost intersect at the developer's computer. By building a tool that measures both phenomena locally, we reveal a 75.6\% reliance rate and a 31x geographic carbon penalty. This sovereignty-preserving telemetry is a necessary step for the Global South to manage the socio-economic transition of widespread AI adoption.

\bibliographystyle{plain}
\begin{thebibliography}{9}

\bibitem{kulveit2025gradual}
Jan Kulveit, et al.
\textit{Gradual Disempowerment: Systemic Existential Risks from Incremental AI Development}.
arXiv:2501.16946, 2025.

\bibitem{metr2025}
METR.
\textit{Measuring the Impact of Early-2025 AI on Experienced Open-Source Developer Productivity}.
arXiv:2507.09089, 2025.

\bibitem{oviedo2025energy}
F. Oviedo, et al.
\textit{Energy Use of AI Inference: Efficiency Pathways and Test-Time Compute}.
Joule; arXiv:2509.20241, 2025.

\bibitem{you2025energy}
J. You / Epoch AI.
\textit{How Much Energy Does ChatGPT Use?}
2025.

\bibitem{greptile2026}
Greptile.
\textit{Rise of the Overnight Agents}.
2026.

\bibitem{hust2024}
Hanoi Univ. of Science \& Technology.
\textit{Vietnam grid emission factor}.
2024.

\bibitem{rte2024}
RTE.
\textit{Annual Electricity Review 2024 (France)}.
2024.

\end{thebibliography}

\end{document}
